\documentclass[10.5pt,a4paper]{article}
\usepackage[top=1.2cm,bottom=1.2cm,left=1.6cm,right=1.6cm]{geometry}
\usepackage[T1]{fontenc}
\usepackage[utf8]{inputenc}
\usepackage{xcolor}
\usepackage{hyperref}

\pagestyle{empty}
\setlength{\parskip}{4.5pt}
\setlength{\parindent}{0pt}

\definecolor{accent}{HTML}{1F4E78}
\hypersetup{colorlinks=true, urlcolor=accent, linkcolor=accent}

\begin{document}

\textbf{\large Prabhavit Gupta}\\
Via Zanetti 10, 30175 Venezia, Italia ~$\vert$~ +39 351 610 7807 ~$\vert$~ \href{mailto:prabhavitg@gmail.com}{prabhavitg@gmail.com}\\
\href{https://www.linkedin.com/in/prabhavit/}{linkedin.com/in/prabhavit} ~$\vert$~ \href{https://github.com/prab-gupta}{github.com/prab-gupta}\\
27 settembre 2026

\textbf{Alla cortese attenzione del Team di Selezione / Forward Deployed Engineering}\\
Accenture Italia\\
Via Delle Industrie, 1 / Sede di Milano\\
20090 Assago (MI) / Milano, Italia\\
Posizione: Forward Deployed Engineer Manager (Rif. LinkedIn: 4444008201)

Gentile Team di Selezione di Accenture Italia,

Desidero sottoporre alla Vostra attenzione la mia candidatura per il ruolo di \textbf{Forward Deployed Engineer Manager} presso Accenture Italia a Milano. Conosco l'impegno strategico con cui Accenture guida la reinvenzione digitale enterprise, colmando il divario tra la sperimentazione dell'intelligenza artificiale e la sua piena industrializzazione tramite team dedicati e partnership globali (tra cui Anthropic, OpenAI e Google). Come ingegnere specializzato nell'implementazione di agenti AI in produzione, integrazioni API complesse e gestione tecnica di prodotto, mi riconosco appieno nella missione del Forward Deployed Engineer: operare direttamente a contatto con i clienti enterprise per guidare risultati operativi concreti, garantendo scalabilita, affidabilita e rapido time-to-value.

Nel corso della mia esperienza maturata su otto realta tecnologiche europee, ho combinato sviluppo software hands-on, coordinamento di team e allineamento con gli stakeholder di business:

\begin{itemize}
  \setlength{\itemsep}{1.5pt}
  \setlength{\parskip}{0pt}
  \item \textbf{Deployment di Agenti AI e Governance di Produzione:} Ho architettato e rilasciato in produzione pipeline multi-agente per l'ingestione e la strutturazione di dati non strutturati (PDF, fogli di calcolo, immagini, web) in schemi JSON validati con revisione human-in-the-loop, redigendo oltre 900 schede operative su FlowSpace. Ho inoltre implementato agenti conversazionali di supporto attivi su database in tempo reale con soglie di confidenza ed escalation automatica. Nel mio lavoro quotidiano adotto \textbf{Claude Code} per accelerare il prototyping, l'orchestrazione degli agenti e la generazione di harness di test automatizzati.
  \item \textbf{Ingegneria di Integrazione e Migrazione Enterprise:} Ho guidato la completa migrazione della piattaforma core di FlowSpace da uno stack legacy a un'architettura Next.js e Laravel (PostgreSQL), mappando i modelli dati e scalando da 2.000 a oltre 4.500 utenti attivi con continuita operativa al 100\%. Ho sviluppato integrazioni REST e webhook bidirezionali (OfficeRnD, Nexudus, CRM, Google APIs) con rigorosa validazione degli schemi, gestione dei rate limit e meccanismi di retry automatici.
  \item \textbf{Leadership Tecnica e Orientamento ai Risultati di Business:} Ho operato come punto di contatto tecnico primario per committenti e partner esterni, traducendo requisiti complessi e ambigui in roadmap ingegneristiche precise. Coordinando pod di sviluppo e interfacciandomi con stakeholder direzionali, garantisco che l'architettura tecnica risponda a metriche di ROI chiare, efficienza dei costi di token e adozione utente.
\end{itemize}

Risiedo a Venezia con rapidi collegamenti ferroviari ad alta velocita per Milano Centrale e Assago, garantendo piena e regolare disponibilita per operare on-site presso i clienti e gli uffici di Accenture. Ho conseguito la laurea in Philosophy, International Studies \& Economics presso l'Universita Ca' Foscari di Venezia, unendo rigore analitico-formale e comprensione dei modelli economici. Sono titolare di un Permesso di Soggiorno UE per Soggiornanti di Lungo Periodo (piena autorizzazione al lavoro in Italia a tempo indeterminato, nessuna sponsorizzazione necessaria), parlo inglese a livello professionale C2 e italiano, e sono disponibile immediatamente.

Sarei entusiasta di approfondire in sede di colloquio come il mio profilo tecnico hands-on, la mia esperienza con agenti AI e la capacita di delivery forward-deployed possano contribuire al successo dei progetti di Accenture Italia e dei suoi clienti enterprise.

Ringraziando per l'attenzione che vorrete dedicarmi, porgo i miei piu cordiali saluti.

Prabhavit Gupta

\vspace{0.4em}
{\footnotesize \textit{Autorizzo il trattamento dei miei dati personali presenti nel presente documento ai sensi del Regolamento UE 2016/679 (GDPR) e del D.Lgs. 196/2003.}\par}

\end{document}